% Oversells «Ish daftari» shabloni (XeLaTeX yoki LuaLaTeX).
% \yozishchiziqlari{3} — yozish chiziqlari, \javobjoyi[20mm] — javob uchun joy,
% \katakjoy[30mm] — katak qogʻoz. Koʻrinish sahifa.sty da («Sahifa» tugmasi).
\documentclass[12pt]{article}
\usepackage[a4paper]{geometry}
\usepackage{fontspec}
\setmainfont{Noto Sans}
\usepackage{babel}
\babelprovide[import=uz-Latn, main]{uzbek}
\def\uzbekhyphenmins{22}
\usepackage{dizayn}
\usepackage{sahifa} % Oversells sahifa dizayni (sahifa.sty)

\begin{document}

\section*{Qoʻshish va ayirish}

\begin{uslub}{topshiriq}
1. Hisoblang va javobni yozing.
\end{uslub}
$7 + 5 = \_\_\_$ \qquad $12 - 4 = \_\_\_$ \qquad $9 + 8 = \_\_\_$
\javobjoyi[15mm]

\begin{uslub}{topshiriq}
2. Masalani yeching.
\end{uslub}
Bogʻda 14 ta olma daraxti bor edi. Yana 6 ta koʻchat ekildi.
Bogʻda nechta daraxt boʻldi?
\begin{uslub}{maslahat}
Maslahat: avval nimani \matnuslubi{urgu}{qoʻshish} kerakligini aniqlang.
\end{uslub}
\yozishchiziqlari{3}

\begin{uslub}{topshiriq}
3. Katakli joyga 5 ta uchburchak chizing.
\end{uslub}
\katakjoy[35mm]

\begin{uslub}{topshiriq}
4. Gapni davom ettiring: «Men eng koʻp yaxshi koʻradigan son …»
\end{uslub}
\yozishchiziqlari{4}

\end{document}
